% This document was generated by an LLM.

\documentclass{essay}

\title{Continuity on the Rational Numbers:\\
What Survives, What Fails, and Why}
\author{}
\date{}

\begin{document}

\maketitle

\section{Introduction}

Continuity is often first encountered as a property of functions
\[
f:\mathbb{R}\to\mathbb{R}.
\]
Because nearly all elementary examples live on the real line, it is easy to
unconsciously associate continuity with many other familiar properties of
\(\mathbb{R}\): the Intermediate Value Theorem, the Extreme Value Theorem,
boundedness on closed intervals, the existence of limits of Cauchy sequences,
and the absence of ``gaps.''

But continuity itself does not require any of these things.

To see this clearly, consider functions
\[
f:\mathbb{Q}\to\mathbb{Q}.
\]
The usual \(\varepsilon\)-\(\delta\) definition makes perfect sense, and many of
the elementary rules of continuity remain valid. Yet some of the most familiar
theorems about continuous real-valued functions fail spectacularly.

This makes \(\mathbb{Q}\) an unusually instructive laboratory. The rational
numbers are dense in \(\mathbb{R}\), so locally they seem to be everywhere.
Nevertheless, they are incomplete and disconnected in a very strong sense.
The resulting examples separate the concept of continuity itself from the
additional structural properties of the real numbers.

The central lesson is:

\[
\boxed{\text{Continuity is relative to the domain. Missing points are invisible to it.}}
\]

\section{Continuity on \(\mathbb{Q}\)}

Let
\[
f:\mathbb{Q}\to\mathbb{Q}.
\]
We say that \(f\) is continuous at \(a\in\mathbb{Q}\) if for every
\(\varepsilon>0\) there exists \(\delta>0\) such that
\[
x\in\mathbb{Q},\qquad |x-a|<\delta
\]
implies
\[
|f(x)-f(a)|<\varepsilon.
\]

This is exactly the usual definition, except that the quantified variable
\(x\) is restricted to rational numbers.

That restriction is not cosmetic. It changes which nearby points the function
is required to control.

Suppose, for example, that something strange happens near the irrational
number \(\sqrt{2}\). If the domain is \(\mathbb{Q}\), there is no requirement
to check continuity \emph{at} \(\sqrt{2}\), because
\[
\sqrt{2}\notin\mathbb{Q}.
\]

The point simply is not part of the space on which the function is defined.

\section{A Pole at a Missing Point}

Consider
\[
f(x)=\frac{1}{x^2-2},
\qquad x\in\mathbb{Q}.
\]

This is a perfectly well-defined function
\[
f:\mathbb{Q}\to\mathbb{Q}.
\]
Indeed, if \(x\in\mathbb{Q}\), then \(x^2\neq 2\), since otherwise
\(x=\pm\sqrt{2}\) would be rational.

As a rational function, \(f\) is continuous at every rational point of its
domain. Therefore \(f\) is continuous everywhere on \(\mathbb{Q}\).

Yet rational values of \(x\) can approach \(\sqrt{2}\) arbitrarily closely,
and along such rational sequences,
\[
\left|\frac{1}{x^2-2}\right|
\]
can become arbitrarily large.

Thus we have a function that is continuous everywhere on its domain and yet
has arbitrarily large values inside the bounded interval
\[
[1,2]\cap\mathbb{Q}.
\]

On \(\mathbb{R}\), such behavior would normally signal a singularity at
\(\sqrt{2}\). On \(\mathbb{Q}\), the singularity sits at a point that does not
exist in the domain.

This already shows that continuity alone does not imply the familiar behavior
of continuous functions on closed real intervals.

\section{A Continuous Step Function}

An even more striking example is
\[
f(x)=
\begin{cases}
0, & x^2<2,\\
1, & x^2>2,
\end{cases}
\qquad x\in\mathbb{Q}.
\]

There is no rational \(x\) satisfying \(x^2=2\), so this defines \(f\) on all
of \(\mathbb{Q}\).

Viewed from \(\mathbb{R}\), the function appears to jump at
\(\pm\sqrt{2}\). But these points are absent from the domain.

Let \(a\in\mathbb{Q}\). Suppose first that \(a^2<2\). Since
\(a\neq \sqrt{2}\) and \(a\neq-\sqrt{2}\), the point \(a\) lies a positive
distance away from the boundary between the regions \(x^2<2\) and \(x^2>2\).
Hence there exists \(\delta>0\) such that every rational \(x\) satisfying
\[
|x-a|<\delta
\]
also satisfies \(x^2<2\). Therefore
\[
f(x)=f(a)=0
\]
throughout that neighborhood.

The same argument applies when \(a^2>2\).

Thus \(f\) is continuous at every rational point.

So a function can be continuous on \(\mathbb{Q}\), take only the values
\(0\) and \(1\), and still change from one value to the other across an
irrational cut.

This is impossible for a continuous function on an interval in
\(\mathbb{R}\).

\section{Failure of the Intermediate Value Theorem}

Take rational numbers \(a,b\) such that
\[
a<\sqrt{2}<b.
\]
For the previous function,
\[
f(a)=0,
\qquad
f(b)=1.
\]

If the Intermediate Value Theorem held, then for every
\[
y\in(0,1)
\]
there would exist some \(c\) between \(a\) and \(b\) satisfying
\[
f(c)=y.
\]

But the image of \(f\) is simply
\[
\{0,1\}.
\]

In particular, there is no \(c\in\mathbb{Q}\) with
\[
f(c)=\frac12.
\]

Thus the Intermediate Value Theorem fails.

The reason is not a defect in continuity. The theorem depends on a separate
property of the real line: connectedness.

A real interval cannot be split into two disjoint nonempty open pieces. By
contrast, the rational numbers can be split at every irrational number.

For example, define
\[
A=\{q\in\mathbb{Q}:q<\sqrt{2}\}
\]
and
\[
B=\{q\in\mathbb{Q}:q>\sqrt{2}\}.
\]
Then
\[
\mathbb{Q}=A\cup B,
\qquad
A\cap B=\varnothing.
\]

Moreover, both \(A\) and \(B\) are open in the subspace topology on
\(\mathbb{Q}\). They are also closed, since each is the complement of the
other.

Such sets are called \emph{clopen}: both closed and open.

The step function above is continuous precisely because it is constant on
these two clopen pieces.

\section{Every Irrational Number Creates a Cut}

The phenomenon is not special to \(\sqrt{2}\).

Let \(\alpha\in\mathbb{R}\setminus\mathbb{Q}\), and define
\[
f_\alpha(q)=
\begin{cases}
0,&q<\alpha,\\
1,&q>\alpha,
\end{cases}
\qquad q\in\mathbb{Q}.
\]

Since \(\alpha\notin\mathbb{Q}\), every rational point lies strictly on one
side or the other. The same argument as before shows that
\[
f_\alpha:\mathbb{Q}\to\mathbb{Q}
\]
is continuous everywhere.

Thus every irrational number determines a nontrivial continuous
\(\{0,1\}\)-valued function on \(\mathbb{Q}\).

This gives a particularly concrete way to see that irrational numbers act as
gaps in the rational line.

From the viewpoint of Dedekind cuts, this is exactly what one should expect.
The irrational number \(\alpha\) corresponds to a partition of the rationals
into those lying below the cut and those lying above it. Before the real
numbers have been constructed, the cut exists as a structural feature of
\(\mathbb{Q}\), even though there is no rational point occupying the boundary.

\section{\(\mathbb{Q}\) Is Totally Disconnected}

The rational numbers are not merely disconnected.

They are \emph{totally disconnected}: the only connected subsets of
\(\mathbb{Q}\) are single points.

To see the idea, take two distinct rational numbers
\[
p<q.
\]
Choose an irrational number \(\alpha\) satisfying
\[
p<\alpha<q.
\]
The sets
\[
(-\infty,\alpha)\cap\mathbb{Q}
\]
and
\[
(\alpha,\infty)\cap\mathbb{Q}
\]
separate \(p\) from \(q\).

Therefore no subset of \(\mathbb{Q}\) containing both \(p\) and \(q\) can
behave like a connected real interval.

This explains why continuous functions on \(\mathbb{Q}\) may jump across
irrational cuts without violating continuity.

The real line and the rational line are both densely ordered, but topologically
they are radically different.

\section{Closed and Bounded Is Not Enough}

Another familiar theorem from elementary analysis says:

\begin{quote}
A continuous real-valued function on a closed bounded interval
\([a,b]\subset\mathbb{R}\) is bounded and attains both a maximum and a minimum.
\end{quote}

This is the Extreme Value Theorem.

Now consider
\[
D=[1,2]\cap\mathbb{Q}
\]
and again let
\[
f(x)=\frac{1}{x^2-2}.
\]

The set \(D\) is closed and bounded as a subset of \(\mathbb{Q}\), and \(f\) is
continuous on \(D\). Nevertheless \(f\) is unbounded.

Indeed, choose rational numbers \(q_n\to\sqrt{2}\). Then
\[
|f(q_n)|
=
\frac{1}{|q_n^2-2|}
\to\infty.
\]

Thus the statement
\[
\text{``closed and bounded implies compact''}
\]
fails in \(\mathbb{Q}\).

The Heine--Borel theorem
\[
K\subset\mathbb{R}
\text{ is compact}
\iff
K\text{ is closed and bounded}
\]
is a theorem specifically about Euclidean space. It is not true in an
arbitrary metric space.

The set
\[
[1,2]\cap\mathbb{Q}
\]
is closed and bounded in \(\mathbb{Q}\), but it is not compact.

The missing irrational limit points prevent compactness.

\section{A Cauchy Sequence With No Limit}

The incompleteness of \(\mathbb{Q}\) appears directly through sequences.

Choose rational numbers \(q_n\) satisfying
\[
q_n\to\sqrt{2}
\]
in \(\mathbb{R}\).

For instance, one may choose decimal truncations
\[
1,\quad
1.4,\quad
1.41,\quad
1.414,\quad
1.4142,\quad\ldots
\]

The sequence \((q_n)\) is Cauchy. For every \(\varepsilon>0\), sufficiently
late terms are arbitrarily close to one another.

But it does not converge in \(\mathbb{Q}\), because its only possible real
limit is \(\sqrt{2}\), which is not rational.

This distinction is subtle but fundamental.

The statement
\[
q_n\to\sqrt{2}
\]
is true when the sequence is considered inside \(\mathbb{R}\).

The same sequence has no limit when considered as a sequence in the metric
space \(\mathbb{Q}\).

Convergence is therefore also relative to the ambient space.

\section{Continuity Does Not Repair Incompleteness}

Consider again
\[
f(x)=\frac{1}{x^2-2}.
\]

Let \(q_n\in\mathbb{Q}\) be a Cauchy sequence with
\[
q_n\to\sqrt{2}
\]
in \(\mathbb{R}\).

Then
\[
f(q_n)=\frac{1}{q_n^2-2}
\]
does not even remain bounded, let alone form a convergent sequence.

This does not contradict continuity, because the Cauchy sequence
\((q_n)\) has no limit in the domain.

Continuity says that if
\[
q_n\to q
\]
for some \(q\in\mathbb{Q}\), then
\[
f(q_n)\to f(q).
\]

It says nothing about Cauchy sequences whose missing limits lie outside the
space.

This is one place where uniform continuity becomes conceptually useful.
Uniformly continuous maps send Cauchy sequences to Cauchy sequences. Ordinary
continuity need not.

Indeed,
\[
f(x)=\frac{1}{x^2-2}
\]
is continuous on
\[
[1,2]\cap\mathbb{Q},
\]
but it is not uniformly continuous there.

If it were uniformly continuous, the Cauchy sequence
\(q_n\to\sqrt{2}\) would have to produce a Cauchy sequence \(f(q_n)\), which it
does not.

\section{A Function Continuous on \(\mathbb{Q}\) but With No Continuous Extension}

A natural question is whether every continuous function
\[
f:\mathbb{Q}\to\mathbb{R}
\]
can be extended to a continuous function
\[
F:\mathbb{R}\to\mathbb{R}.
\]

The answer is no.

The step function
\[
f(q)=
\begin{cases}
0,&q<\sqrt{2},\\
1,&q>\sqrt{2}
\end{cases}
\]
provides a counterexample.

Suppose a continuous extension \(F:\mathbb{R}\to\mathbb{R}\) existed.

Choose rational sequences
\[
a_n<\sqrt{2},
\qquad
a_n\to\sqrt{2},
\]
and
\[
b_n>\sqrt{2},
\qquad
b_n\to\sqrt{2}.
\]

Then
\[
F(a_n)=0
\]
for every \(n\), so continuity of \(F\) at \(\sqrt{2}\) would force
\[
F(\sqrt{2})=0.
\]

But also
\[
F(b_n)=1
\]
for every \(n\), so continuity would force
\[
F(\sqrt{2})=1.
\]

This is impossible.

Thus continuity on a dense subset does not automatically determine a
continuous extension to its completion.

Additional conditions are required. Uniform continuity is one important
condition that does allow extension from a dense subset into a complete target
space.

\section{Density Is Not Completeness}

One source of confusion is that the rational numbers are dense in the real
numbers.

Between any two distinct real numbers there is a rational number.

This may create the impression that \(\mathbb{Q}\) is somehow ``almost all''
of \(\mathbb{R}\), at least for questions involving limits.

But density and completeness express entirely different ideas.

Density says, roughly:
\[
\text{there are points of the set arbitrarily close to every real point.}
\]

Completeness says:
\[
\text{Cauchy processes that should converge actually have their limits
inside the space.}
\]

The sequence
\[
1,\ 1.4,\ 1.41,\ 1.414,\ldots
\]
illustrates both properties simultaneously.

Its terms are rational, showing that rationals can approximate \(\sqrt{2}\)
arbitrarily well.

Yet its limit is missing from \(\mathbb{Q}\).

Density provides arbitrarily good approximations.

Completeness provides the point being approximated.

\section{A Strange Contrast: Locally Rich, Globally Full of Gaps}

The rational numbers have an unusual combination of properties.

They are dense:
every real interval contains infinitely many rational numbers.

They have no isolated points:
every rational number has other rational numbers arbitrarily close to it.

Yet they are totally disconnected:
no two distinct rational numbers belong to a connected subset of
\(\mathbb{Q}\).

This is one of the reasons \(\mathbb{Q}\) is so useful pedagogically.

Local inspection suggests an extremely rich space. Around every rational
point, there are infinitely many other rational points.

But globally, irrational cuts split the space everywhere.

This shows that ``having many nearby points'' and ``being connected'' are
completely different notions.

\section{Another Paradoxical Example: A Continuous Function Taking Only
Rationally Discrete Values}

Let \(\alpha\) be irrational and define
\[
f(q)=
\begin{cases}
17,&q<\alpha,\\
-5,&q>\alpha.
\end{cases}
\]

This function is continuous on \(\mathbb{Q}\).

Nothing special about \(0\) and \(1\) was used. The values on the two sides
can be separated by an arbitrarily large jump.

Thus for every \(M>0\), there exists a continuous function on
\(\mathbb{Q}\) having a jump of size larger than \(M\) across an irrational
cut.

The phrase ``jump discontinuity'' therefore has to be used carefully.
There is visually a jump when the rational graph is embedded in the real
plane, but there is no discontinuity at any point of the domain.

A discontinuity must occur \emph{at a point where the function is being asked
to be continuous}.

\section{Continuity Is Intrinsic to the Chosen Space}

Suppose a function is defined by
\[
f(q)=
\begin{cases}
0,&q<\sqrt{2},\\
1,&q>\sqrt{2},
\end{cases}
\qquad q\in\mathbb{Q}.
\]

If we regard its domain as \(\mathbb{Q}\), then it is continuous.

If we enlarge the domain by adding the missing point \(\sqrt{2}\), then no
choice of the value \(f(\sqrt{2})\) makes the resulting function continuous.

If we enlarge the domain all the way to \(\mathbb{R}\), we obtain the familiar
step-function phenomenon.

So continuity is not merely a property of a formula.

It is a property of a function together with its domain and codomain.

The formula alone does not determine the topological question.

This is a useful general principle well beyond this example. In analysis and
topology one must always ask:

\begin{itemize}
    \item What is the domain?
    \item What topology or metric is being used?
    \item Which points belong to the space?
    \item Which limiting processes are required to remain inside the space?
\end{itemize}

\section{What Actually Depends on Completeness?}

It is worth separating several ideas that are often first learned together.

The definition of continuity does not require completeness.

The sequential characterization of continuity also does not require
completeness.

The algebraic rules
\[
f,g\text{ continuous}
\implies
f+g,\ fg,\ f\circ g\text{ continuous}
\]
remain valid in the appropriate metric-space setting.

By contrast, completeness enters when one wants to guarantee that certain
limiting processes actually produce points inside the space.

For example:
\begin{itemize}
    \item Cauchy sequences need not converge in \(\mathbb{Q}\).
    \item Bounded subsets of \(\mathbb{Q}\) need not have rational suprema.
    \item Closed and bounded subsets of \(\mathbb{Q}\) need not be compact.
\end{itemize}

Connectedness is responsible for a different family of results:
\begin{itemize}
    \item the Intermediate Value Theorem;
    \item the fact that continuous images of intervals are intervals;
    \item the impossibility of nonconstant continuous maps from an interval
    into a discrete two-point set.
\end{itemize}

Compactness is responsible for yet another family:
\begin{itemize}
    \item continuous functions on compact sets are bounded;
    \item continuous functions on compact sets attain maxima and minima;
    \item continuous functions on compact sets are uniformly continuous.
\end{itemize}

Thus several theorems that appear under the broad heading of ``continuity''
are actually consequences of continuity combined with additional structure.

\section{A Useful Diagnostic Table}

\[
\begin{array}{c|c}
\text{Statement} & \text{On }\mathbb{Q} \\ \hline
\varepsilon\text{-}\delta\text{ continuity makes sense}
& \text{Yes}\\
\text{Sequential characterization of continuity}
& \text{Yes}\\
\text{Sums and compositions of continuous functions are continuous}
& \text{Yes}\\
\text{Every Cauchy sequence converges}
& \text{No}\\
\text{Every bounded set has a supremum in the space}
& \text{No}\\
\text{Intermediate Value Theorem}
& \text{No}\\
\text{Closed and bounded implies compact}
& \text{No}\\
\text{Continuous on closed bounded set implies bounded}
& \text{No}\\
\text{Continuous on closed bounded set implies uniformly continuous}
& \text{No}
\end{array}
\]

The table makes an important methodological point visible.

When a theorem about continuous real functions fails on \(\mathbb{Q}\), one
should not immediately conclude that ``continuity behaves differently.''

Instead ask:

\begin{quote}
Which additional property of \(\mathbb{R}\) was the theorem secretly using?
\end{quote}

Very often the answer is connectedness, compactness, or completeness.

\section{Dedekind Cuts Revisited}

These examples also illuminate the motivation for constructing the real
numbers.

Suppose we remain entirely inside \(\mathbb{Q}\).

The sets
\[
A=\{q\in\mathbb{Q}:q^2<2,\ q>0\}
\]
and the corresponding rational numbers lying above the cut reveal a structural
gap. Rational numbers approach the boundary arbitrarily closely from both
sides, but no rational number occupies it.

A Dedekind construction does not merely declare that a mysterious irrational
number must exist somewhere outside \(\mathbb{Q}\).

Instead, the cut itself is used as the new real number.

In this sense, the continuous step function
\[
f(q)=
\begin{cases}
0,&q<\sqrt{2},\\
1,&q>\sqrt{2}
\end{cases}
\]
is not merely an amusing counterexample.

It is a direct topological manifestation of the same gap that the Dedekind
construction is designed to fill.

Once the missing point is inserted, the two rational sides acquire an actual
boundary point. Continuity across that completed space becomes more
restrictive.

\section{Conclusion}

Functions on \(\mathbb{Q}\) reveal that continuity is a much weaker and more
local concept than its first appearance in elementary calculus suggests.

A function may be continuous everywhere on \(\mathbb{Q}\) while:

\begin{itemize}
    \item becoming arbitrarily large near a missing irrational point;
    \item jumping from \(0\) to \(1\) across an irrational cut;
    \item violating the Intermediate Value Theorem;
    \item being unbounded on a closed bounded subset of \(\mathbb{Q}\);
    \item failing to be uniformly continuous on such a set;
    \item admitting no continuous extension to \(\mathbb{R}\).
\end{itemize}

None of this contradicts the definition of continuity.

The apparent paradoxes arise because theorems familiar from real analysis
usually combine continuity with deeper properties of the real numbers.

The rational numbers are dense but incomplete.

They contain infinitely many points in every interval but are totally
disconnected.

They permit arbitrarily accurate approximation to irrational limits while
failing to contain those limits.

For this reason, \(\mathbb{Q}\) provides an unusually clear way to disentangle
four concepts that are easy to blur together:

\[
\boxed{
\text{continuity},
\qquad
\text{connectedness},
\qquad
\text{compactness},
\qquad
\text{completeness}.
}
\]

Understanding exactly which theorem depends on which property is one of the
central conceptual transitions from calculus to real analysis.

\end{document}
